% LLNCS report for CENG 467 Final Project
% Faithfulness Under Pressure: Context and History Pruning for RAG
%
\documentclass[runningheads]{llncs}
%
\usepackage[T1]{fontenc}
\usepackage{graphicx}
\usepackage{booktabs}
\usepackage{amsmath}
\usepackage{amssymb}
\usepackage{array}
\usepackage{tikz}
\usetikzlibrary{arrows.meta}
\usepackage{caption}
\usepackage{url}
%
% Figure search paths (relative to this .tex file)
\graphicspath{%
  {../../simulations/cfc_proof/figures/}%
  {../../outputs/figures/}%
  {../../experiments/results/notebook5/}%
}
%
\usepackage{color}
\renewcommand\UrlFont{\color{blue}\rmfamily}
\urlstyle{rm}
%
\begin{document}
%
\title{Faithfulness Under Pressure: Context and\\ History Pruning Strategies for RAG Faithfulness\\ in Small Language Models}
%
\titlerunning{Faithfulness Under Pressure}
%
\author{Mert G\"{u}den\inst{1} \and Berkay Fehmi Tekin\inst{1}}
%
\authorrunning{M. G\"{u}den and B. F. Tekin}
%
\institute{CENG 467 -- Natural Language Understanding and Generation\\
\email{\{300201013, 300201090\}}\\
\url{https://github.com/Mrtuzy/CENG467_Final}}
%
\maketitle
%
\begin{abstract}
Retrieval-augmented generation (RAG) grounds language model outputs in
retrieved documents, but small language models (SLMs) remain vulnerable to
the noisy, lengthy contexts that retrieval produces. We study whether
\emph{context pruning} can make SLM answers more faithful to the retrieved
evidence. We implement a modular pruning harness with seven strategies
spanning no pruning, naive truncation, RECOMP-style extraction, LLMLingua-2
token compression, Provence sentence pruning, and history-aware pruning, all
held under identical retrieval and generation conditions. Faithfulness is
scored with a decompose-then-verify LLM-as-Judge. On a HotpotQA distractor
pilot, real LLMLingua-2 attains the best judge faithfulness (0.711) but at
near-zero verbatim supporting-fact coverage, exposing a coverage--noise
tension that single metrics hide. Motivated by the conversational case, we
further design an entropy-driven adaptive pruner built on Closed-form
Continuous-time (CfC) networks, in which per-turn surprisal sets a continuous
integration time step. A seven-simulation verification suite fixes every
hyperparameter analytically and produces a falsifiable go/no-go test; on real
QReCC teacher labels the surprisal--importance correlation is
$-0.074$, below the $\rho^{*}=0.1$ threshold, and the trained CfC fails to
generalise. We report this as an honest negative result and discuss its
implications for adaptive pruning.
\keywords{Retrieval-augmented generation \and Context pruning \and
Faithfulness \and LLM-as-Judge \and Small language models}
\end{abstract}
%
\section{Introduction and Problem Definition}

Retrieval-augmented generation (RAG)~\cite{lewis2020rag} grounds the output of
a language model in documents fetched at inference time, improving performance
on knowledge-intensive tasks and reducing
hallucination~\cite{shuster2021retrieval}. The promise of RAG, however, comes
with a cost: the retriever returns long passages that mix genuine evidence with
near-misses, stale entities, and distractors. Small language models (SLMs),
which are attractive for cost and latency, are especially sensitive to such
noise, and long-context models are known to underuse evidence buried in the
middle of a prompt~\cite{liu2024lost}.

This project asks a single question: \emph{when an SLM receives long, noisy
retrieved context, can pruning make its answer more faithful to the evidence?}
\emph{Faithfulness}---the degree to which a generated answer is grounded in the
retrieved context---is a central quality dimension for RAG, yet it is hard to
measure automatically and is precisely what degrades under noisy context.

\paragraph{Task.}
We work on open-domain question answering with RAG. Our primary experiments use
HotpotQA~\cite{yang2018hotpotqa} in the distractor setting, a multi-hop QA
benchmark where each question is paired with ten passages: two supporting and
eight distractors drawn from Wikipedia. The pipeline retrieves the top
passages with BM25, optionally applies a pruning strategy to reduce length and
noise, generates a short answer with an instruction-tuned SLM, and scores the
answer's faithfulness with an LLM-as-Judge that decomposes the answer into
atomic claims and verifies each against the context. The input is a question
(plus optional dialogue history); the output is a short answer and its
faithfulness score. The task is therefore open-domain QA coupled to a
reference-free \emph{evaluation} sub-task.

\paragraph{Technical challenge.}
The difficulty is threefold. First, faithfulness is not directly measurable
without expensive human annotation; we approximate it with a calibrated
LLM-as-Judge~\cite{es2023ragas} anchored by Cohen's $\kappa$ against human
labels. Second, pruning introduces a \emph{coverage--noise trade-off}:
aggressive compression can remove the very sentences that support the correct
answer. Third, SLMs hallucinate more under long or noisy contexts, making
pruning quality decisive. The seven strategies we evaluate span from no pruning
to combined semantic and history-aware compression, letting us isolate the
effect of each design choice.

\paragraph{Topic evolution.}
The work began on multi-hop distractor QA (HotpotQA) and progressively shifted
toward the \emph{conversational} case (QReCC~\cite{anantha2021qrecc}), where the
central problem is deciding which past turns and retrieved passages still matter
for the current question. HotpotQA runs are retained as implementation sanity
checks; the conversational direction motivates the adaptive CfC pruner of
Section~\ref{sec:cfc}.

\section{Dataset and Benchmark Status}

We use HotpotQA (distractor) loaded via the HuggingFace \texttt{datasets}
library (\texttt{load\_dataset("hotpot\_qa", "distractor")}), with roughly
90{,}447 training and 7{,}405 validation samples. We fine-tune no generator;
all experimental samples come from the validation split. Each sample provides a
question, a short gold answer (a span or yes/no value), ten titled passages,
and supporting-fact labels marking the relevant sentences. During preprocessing
we flatten each passage's sentences into a single string, apply whitespace
tokenization for BM25 indexing, and truncate passages to a configurable token
budget. The loader (\texttt{src/data/loader.py}) writes cleaned JSONL records,
which are ready for the retrieval and pruning stages.

The conversational benchmark, QReCC~\cite{anantha2021qrecc}, contains about
14K conversations and 81K question--answer pairs over a 54M-passage web
collection, with each sample exposing the original question, conversational
context, a human rewrite, the answer, and answer provenance. It is the intended
final benchmark for history-aware pruning; the CfC study (Section~\ref{sec:cfc})
already uses QReCC-derived teacher labels.

\begin{table}[t]
\centering
\caption{HotpotQA dataset summary (distractor setting).}\label{tab:dataset}
\begin{tabular}{ll}
\toprule
\textbf{Property} & \textbf{Value}\\
\midrule
Full name & HotpotQA (distractor setting)\\
Source & HuggingFace \texttt{hotpot\_qa}\\
Training / validation samples & $\sim$90{,}447 / $\sim$7{,}405\\
Samples used (experiments / calibration) & 200 / 50\\
Passages per question & 10 (2 supporting + 8 distractors)\\
Answer type & Short span / yes--no\\
Task type & Multi-hop open-domain QA\\
Status & Loader implemented, JSONL prepared\\
\bottomrule
\end{tabular}
\end{table}

The main challenges HotpotQA introduces are: (1)~the eight distractors directly
inflate context length and inject factual noise; (2)~multi-hop questions need
evidence from at least two passages, so aggressive pruning risks discarding a
required supporting passage; and (3)~there is no dedicated retrieval split, so
we evaluate end-to-end retrieval alongside pruning on the same validation
samples.

\section{Literature Review}

\paragraph{RAG and faithfulness.}
RAG~\cite{lewis2020rag} grounds generation in retrieved documents and improves
knowledge-intensive tasks; \cite{shuster2021retrieval} shows retrieval reduces
conversational hallucination. FActScore~\cite{min2023factscore} decomposes
long-form text into atomic facts and verifies each against a knowledge source,
directly motivating our claim-level faithfulness metric.
\cite{liu2024lost} documents the ``lost in the middle'' effect, motivating
pruning as a way to surface evidence.

\paragraph{Context compression.}
LLMLingua~\cite{jiang2023llmlingua} scores tokens with a small proxy LM and
drops low-importance ones; LLMLingua-2~\cite{pan2024llmlingua2} reframes
compression as token classification with data distillation, reporting
$2$--$5\times$ compression with $1.6$--$2.9\times$ latency speedups.
Provence~\cite{yoon2024provence} treats sentence pruning as sequence labelling
for RAG at negligible added cost. RECOMP~\cite{xu2023recomp} frames compression
as extractive or abstractive summarization and reports compression rates as low
as 6\% with minimal performance loss. These four works define our baseline and
advanced pruners.

\paragraph{LLM-as-Judge.}
\cite{zheng2023mtbench} shows strong LLMs correlate well with human rankings
despite positional and verbosity biases. RAGAS~\cite{es2023ragas} provides a
reference-free RAG evaluation that checks whether each claim in the answer is
entailed by the context---the decompose-then-verify paradigm we adopt, with
open-weight rather than proprietary models, calibrated by Cohen's $\kappa$.

\paragraph{Adaptive computation.}
For the conversational extension we draw on Closed-form Continuous-time (CfC)
networks~\cite{hasani2022cfc}, which capture continuous temporal dynamics
without a heavy ODE solver, and on surprisal theory from
psycholinguistics~\cite{levy2008,smith2013}, where reading time scales
logarithmically with word surprisal $-\log P(w_j\mid w_{<j})$. We test whether
surprisal is a useful proxy for per-turn importance.

\section{Baseline and Advanced Pruning Strategies}

We evaluate seven strategies (Table~\ref{tab:strategies}) organized as three
baselines and four advanced methods, all sharing a common \texttt{BasePruner}
interface and registered in a central \texttt{PRUNER\_REGISTRY} so that
retrieval and generation are held fixed and only the retained context changes.

\begin{table}[t]
\centering
\caption{Pruning strategies under evaluation. All implemented and complete.}\label{tab:strategies}
\begin{tabular}{lll}
\toprule
\textbf{ID} & \textbf{Strategy} & \textbf{Type}\\
\midrule
B1 & No pruning (full context)        & Baseline (coverage upper bound)\\
B2 & Naive truncation (512 tokens)    & Baseline (production fallback)\\
B3 & RECOMP-style extractive          & Prior-work extractive baseline\\
S1 & LLMLingua-2                      & Token-level compression\\
S2 & Provence                        & Sentence-level pruning\\
S3 & Semantic history pruning         & History-aware pruning\\
S4 & Provence + history (combined)    & Combined\\
\bottomrule
\end{tabular}
\end{table}

B1 is the coverage upper bound: maximum evidence but maximum noise. B2 is the
standard engineering fallback under a hard token budget. B3 selects
query-relevant sentences by lexical (Jaccard) overlap, bridging truncation and
learned compression. S1 and S2 are the state of the art in token- and
sentence-level compression. S3 removes passages that overlap with information
already surfaced in prior turns. S4 chains Provence then history pruning to
combine semantic filtering with redundancy removal.

\section{Experimental Results}

All experiments were reproduced in Google Colab with GPU runtimes.
Mistral-7B-Instruct-v0.2 serves as both generator and judge (shared when
possible). Notebook~2 is a broad seven-pruner sanity run over 50 HotpotQA
validation samples; Notebook~5 is a corrected gap-closing run over a held-out
50-sample slice (validation indices 600--649) that fixes three weaknesses:
genuine LLMLingua-2 execution, non-empty two-turn history pruning, and
qualitative error exports.

\subsection{Baseline Sanity Run}

Table~\ref{tab:baseline} confirms the expected pattern: RECOMP improves
faithfulness over hard truncation, while Provence's default threshold is very
aggressive for multi-hop HotpotQA. The S1 row here is a historical artifact:
LLMLingua-2 fell back to the RECOMP implementation, so S1 and B3 are identical;
S3/S4 use empty history and are degenerate single-turn cases.

\begin{table}[t]
\centering
\caption{Notebook~2 baseline results --- 50 HotpotQA validation samples,
Mistral-7B-Instruct-v0.2. Comp.\ = compression ratio (kept/input);
Lat.\ = latency.}\label{tab:baseline}
\begin{tabular}{lccccr}
\toprule
\textbf{Strategy} & \textbf{Faith.} & \textbf{EM} & \textbf{F1} & \textbf{Comp.} & \textbf{Lat.\ (s)}\\
\midrule
B1: No pruning            & 0.718 & 0.040 & 0.148 & 1.000 & 6.76\\
B2: Naive truncation      & 0.621 & 0.040 & 0.139 & 0.918 & 4.95\\
B3: RECOMP                & \textbf{0.750} & 0.040 & 0.141 & 0.911 & 5.38\\
S1: LLMLingua-2 (fallback)& \textbf{0.750} & 0.040 & 0.141 & 0.911 & 5.38\\
S2: Provence              & 0.657 & 0.000 & 0.057 & \textbf{0.072} & \textbf{4.49}\\
S3: History (empty)       & 0.718 & 0.040 & 0.148 & 1.000 & 6.34\\
S4: Combined (empty)      & 0.657 & 0.000 & 0.057 & \textbf{0.072} & \textbf{4.48}\\
\bottomrule
\end{tabular}
\end{table}

\subsection{Gap-Closing Run}

Notebook~5 re-ran the problematic strategies under corrected conditions:
LLMLingua-2 with the \texttt{llmlingua} package and \texttt{use\_llmlingua2=True}
(zero fallbacks recorded), and S3/S4 with synthetic two-turn history. We add
supporting-fact \emph{coverage} (fraction of gold supporting sentences retained)
and \emph{noise} ratio to make pruning failures visible beyond faithfulness
alone (Table~\ref{tab:gap}, Fig.~\ref{fig:tradeoff}).

\begin{table}[t]
\centering
\caption{Notebook~5 gap-closing results --- 50 held-out HotpotQA samples
(indices 600--649). Cov.\ = supporting-fact coverage.}\label{tab:gap}
\begin{tabular}{lcccccc}
\toprule
\textbf{Strategy} & \textbf{Faith.} & \textbf{EM} & \textbf{F1} & \textbf{Comp.} & \textbf{Cov.} & \textbf{Noise}\\
\midrule
B3: RECOMP              & 0.648 & 0.080 & 0.195 & 0.962 & \textbf{0.760} & 0.624\\
S1: LLMLingua-2 (real)  & \textbf{0.711} & 0.040 & 0.167 & 0.457 & 0.000 & 0.792\\
S3: History (two-turn)  & 0.533 & \textbf{0.100} & \textbf{0.214} & 0.962 & 0.627 & 0.641\\
S4: Combined (two-turn) & 0.530 & 0.040 & 0.112 & \textbf{0.081} & 0.222 & \textbf{0.284}\\
\bottomrule
\end{tabular}
\end{table}

\begin{figure}[t]
\centering
\includegraphics[width=0.72\textwidth]{faithfulness_vs_compression.png}
\caption{Faithfulness versus compression on the HotpotQA gap-closing run. The
desirable region is high faithfulness with a smaller context. LLMLingua-2
reaches the best faithfulness at much higher compression, while the combined
pruner is the most compact but discards too much supporting
evidence.}\label{fig:tradeoff}
\end{figure}

Corrected LLMLingua-2 is no longer identical to RECOMP: it compresses far more
aggressively (0.457 vs.\ 0.962) and obtains higher judge faithfulness (0.711
vs.\ 0.648), yet its verbatim supporting-fact coverage is 0.000. LLMLingua-2
often rewrites or fragments evidence in a way that can still support the answer
but no longer preserves gold sentences verbatim---so coverage and faithfulness
measure different things. Two-turn history pruning changes behaviour relative to
the empty-history case but here lowers faithfulness to 0.533, suggesting the
novelty-only redundancy score discards useful repeated evidence. The combined
setting is the most compact (0.081) and lowest-noise (0.284), but its low
coverage (0.222) confirms the Provence threshold is too aggressive for multi-hop
QA. The broad lesson is that faithfulness, EM/F1, and coverage capture distinct
failure modes, so no single metric should be treated as ground truth.

\section{Error Analysis and Generation Quality}

Notebook~5 exports concrete error candidates
(\texttt{error\_analysis\_candidates.jsonl}). Four recurring failure modes
emerge (Table~\ref{tab:errors}).

\begin{table}[t]
\centering
\caption{Recurring error categories with representative examples.}\label{tab:errors}
\begin{tabular}{p{0.30\textwidth}p{0.62\textwidth}}
\toprule
\textbf{Error type} & \textbf{Example pattern}\\
\midrule
Pruning removes evidence &
``Who wrote the song \emph{Arizona}?'' loses the sentence naming
\emph{Kenny Young}; LLMLingua-2 answers ``I don't know,'' coverage 0.000.\\
Verbose-but-correct under-credited &
For ``Are both Parodia and Thalictrum flowering plants?'' S4 answers
correctly with explanation: faithfulness 1.0, coverage 1.0, but EM 0.0
(gold is ``yes'').\\
Judge false negative &
For Watercliffe Meadow, the answer ``Political correctness'' gets
EM/F1 = 1.0, coverage 1.0, but judge faithfulness 0.0.\\
Entity near-miss &
Generator outputs ``T.\ M.\ Howard'' instead of ``T.\ R.\ M.\ Howard'';
F1 0.857 but EM 0.0 and faithfulness 0.0.\\
\bottomrule
\end{tabular}
\end{table}

The strongest pattern is that faithfulness, EM/F1, and coverage disagree by
construction: LLMLingua-2 has the best faithfulness yet zero retained-sentence
coverage, while some exactly correct short answers receive low judge scores.
This makes human calibration of the judge a prerequisite before treating
faithfulness as a standalone target, and motivates a semantic rather than
verbatim coverage metric.

\section{Adaptive Pruning: An Entropy-Driven CfC Extension}\label{sec:cfc}

Fixed thresholds cannot tell whether a conversation turn is easy, noisy,
redundant, or load-bearing for the current question. We explore an adaptive
controller that allocates more ``processing time'' to surprising turns,
borrowing from surprisal theory: reading time scales with surprisal
$-\log P(w_j\mid w_{<j})$~\cite{levy2008,smith2013}. We do \emph{not} claim
surprisal equals importance; we test whether it is a useful proxy.

\paragraph{Mechanism.}
For utterance $u_i$ we compute mean token surprisal with a DistilGPT-2 proxy,
$S(u_i) = -\tfrac{1}{N}\sum_{j=1}^{N}\log P(w_j\mid w_{<j})$, and map it to a
continuous CfC time step (rank-normalised for outlier robustness):
\begin{equation}
\Delta t_i = \Delta t_{\min} + \beta\,\frac{\mathrm{rank}(S_i)}{T},
\qquad
\varphi_i = 1 - e^{-\lambda\,\Delta t_i}.
\end{equation}
The CfC integrates over $\Delta t_i$ each turn; the commitment fraction
$\varphi_i$ measures how strongly the hidden state locks on, so a larger
$\Delta t_i$ leaves a stronger trace and yields a higher importance score
$\hat p_i = \sigma(Wh_i + b)$. Turns with $\hat p_i < \tau$ are pruned.

\paragraph{Simulate before training.}
Because end-to-end tuning with a 7B teacher is expensive and the proxy may carry
no signal at all, we ran seven model-free simulations, each isolating one design
question and returning a config value or a falsifiable go/no-go threshold
(Table~\ref{tab:sims}). The recurring metric is \emph{importance recovery}: the
Spearman correlation between a method's per-turn score and the latent true
importance.

\begin{table}[t]
\centering
\caption{Simulation suite: each item fixes a parameter or a go/no-go test.}\label{tab:sims}
\begin{tabular}{cp{0.42\textwidth}p{0.40\textwidth}}
\toprule
\textbf{\#} & \textbf{Question} & \textbf{Result}\\
\midrule
S0 & Does the commitment lever exist in the real \texttt{ncps} CfC? & $\rho=+0.944$; analytic error $\approx 4\times10^{-10}$\\
S1 & Is surprisal$\to$commitment a tautology? & No --- closed-form; lever saturates at high $S$\\
S2 & Which $\beta^{*},\Delta t_{\min}^{*}$ maximise recovery? & $\beta^{*}=1.0$, $\Delta t_{\min}^{*}=0.05$\\
S5 & Which $\Delta t$ mapping survives outliers? & Rank normalisation most robust\\
S4 & What does leave-one-out KL measure? & Marginal answer-necessity = correct target\\
S3 & Does surprisal proxy importance well enough? & Go/no-go: $\rho^{*}\!\approx\!0.1$; $\rho\!\geq\!0.5$ beats cosine\\
S6 & What fidelity--compression is reachable? & 40\% compression at 90\% fidelity\\
\bottomrule
\end{tabular}
\end{table}

Two simulations decide the project. S2 (Fig.~\ref{fig:s2}) shows recovery peaks
at $\beta^{*}=1.0$, $\Delta t_{\min}^{*}=0.05$, and---critically---at $\beta=0$
(uniform $\Delta t$) recovery collapses to $\approx 0$: the entropy term, not the
network, carries the signal. S3 (Fig.~\ref{fig:s3}) is the honest go/no-go test:
entropy-$\Delta t$ only beats uniform once
$\rho(\text{surprisal},\text{importance})>\rho^{*}\approx0.1$ and only beats a
cosine baseline once $\rho\geq0.5$. The action is to measure $\rho$ on real
QReCC \emph{before} training.

\begin{figure}[t]
\centering
\includegraphics[width=0.92\textwidth]{s2_beta_dtmin_design.png}
\caption{S2 --- choosing $\beta$ and $\Delta t_{\min}$ from data. Left: recovery
over the $(\beta,\Delta t_{\min})$ grid peaks at the starred point. Right: at
$\beta=0$ recovery is near zero, confirming the entropy term carries the
signal.}\label{fig:s2}
\end{figure}

\begin{figure}[t]
\centering
\begin{minipage}[c]{0.49\textwidth}
\centering
\includegraphics[width=\linewidth]{s3_proxy_quality_limit.png}
\caption{S3 --- recovery rises with proxy quality $\rho$; go/no-go at
$\rho^{*}\approx0.1$, cosine-beating at $\rho\geq0.5$.}\label{fig:s3}
\end{minipage}\hfill
\begin{minipage}[c]{0.49\textwidth}
\centering
\includegraphics[width=\linewidth]{proxy_quality_gonogo.png}
\caption{Real QReCC: mean surprisal--teacher Spearman is $-0.074$, below the
$\rho^{*}=0.1$ threshold.}\label{fig:gonogo}
\end{minipage}
\end{figure}

\paragraph{Real-data result: an honest negative.}
On QReCC teacher labels (answer-conditioned leave-one-out $\Delta$NLL of the
gold answer, scored with Qwen2.5-1.5B-Instruct), the mean surprisal--importance
Spearman is $-0.074$ (Fig.~\ref{fig:gonogo})---below the go/no-go threshold and
far below the cosine-beating target. We trained the CfC anyway as a stress test
(SBERT turn embeddings + rank-normalised $\Delta t_i$, MSE to normalised teacher
importance). It did not generalise (Fig.~\ref{fig:train},
Fig.~\ref{fig:valmetrics}): train loss fell from 0.1293 to 0.1201 while
validation loss barely moved (best 0.1251 at epoch~6, then drifting to 0.1257);
best validation F1 was only 0.050 at $\tau=0.5$, recall peaked at 0.027, and MAE
stayed flat near 0.300. The model predicts very few important turns---high
precision (0.397) only because it abstains.

\begin{figure}[t]
\centering
\includegraphics[width=0.62\textwidth]{train_val_loss.png}
\caption{CfC training and validation loss. Training improves while validation
is flat then worsens after epoch~6 --- memorisation, not importance
learning.}\label{fig:train}
\end{figure}

\begin{figure}[t]
\centering
\begin{minipage}[c]{0.49\textwidth}
\centering
\includegraphics[width=\linewidth]{val_f1_curve.png}
\end{minipage}\hfill
\begin{minipage}[c]{0.49\textwidth}
\centering
\includegraphics[width=\linewidth]{val_mae_curve.png}
\end{minipage}
\caption{CfC validation F1 (left) and MAE (right). F1 peaks at 0.050; MAE stays
near 0.300, indicating uncalibrated per-turn scores.}\label{fig:valmetrics}
\end{figure}

The CfC mechanism is mathematically valid (S0 confirms the commitment lever at
$\rho=+0.944$), but this experiment \emph{rejects} the stronger claim that
DistilGPT-2 surprisal is a reliable importance proxy for QReCC history turns.
The simulation suite paid off: it predicted the failure cheaply and told us
exactly which assumption broke.

\section{Ablation and Prompt-Sensitivity Plan}\label{sec:ablation}

We plan a targeted ablation set to attribute faithfulness gains to the
compression strategy rather than incidental factors: (1)~\emph{judge prompt
format}, zero-shot vs.\ 2-shot for both claim decomposition and verification;
(2)~\emph{verification response}, strict binary YES/NO vs.\ a soft confidence
thresholded at 0.5; (3)~\emph{retrieval depth}, $k=3$ vs.\ $k=5$;
(4)~\emph{truncation budget}, 256 vs.\ 512 tokens for B2/B3; and
(5)~\emph{history context}, empty (single-turn) vs.\ two-turn history for S3/S4.
For the CfC extension the ablation axes are the proxy channel (raw surprisal vs.\
question-conditioned similarity vs.\ LOO~$\Delta$NLL features) and the pruning
rule (fixed $\tau=0.5$ vs.\ ranking loss / top-$k$).

\section{Ethical Considerations and Bias}

The judge is Mistral-7B-Instruct, which may inherit positional bias (favouring
early claims), verbosity bias (favouring longer answers), and cultural skews
from its training data. HotpotQA passages come from English Wikipedia, which
underrepresents non-Western entities, so faithfulness estimates may be more
reliable for densely covered topics. Crucially, faithfulness to context is not
truthfulness about the world: a model that faithfully reproduces a false
retrieved claim scores high while being wrong. All models are open-weight and
used solely for academic research; no proprietary APIs and no personal data are
involved. Results should not be extrapolated to high-stakes QA without
additional human evaluation.

\section{Reproducibility Status}

The repository (\url{https://github.com/Mrtuzy/CENG467_Final}) is a modular
Python package with separated data, retrieval, pruning, generation, evaluation,
and judge components, plus a simulation suite for the CfC study. The primary
execution interface is Google Colab:
\texttt{01\_data\_exploration.ipynb} (dataset checks),
\texttt{02\_baseline\_results.ipynb} (baseline sanity),
\texttt{05\_final\_colab\_runs.ipynb} (LLMLingua-2, two-turn history, judge
calibration scaffold, error exports), and the LNN/CfC notebooks as exploratory
artifacts. All seven pruners are implemented behind the shared registry, and
result JSONL files, an aggregate summary, and error candidates are written under
\texttt{experiments/results/}.

\section{Current Challenges and Next Steps}

Three challenges remain. First, aggressive compression still trades coverage for
faithfulness: real LLMLingua-2 reaches the best faithfulness (0.711) at 0.000
verbatim coverage, so we need a semantic coverage metric. Second, Provence's
default threshold is too aggressive for multi-hop QA (combined keeps 8.1\% of
tokens at 0.222 coverage); threshold tuning toward a 0.4--0.6 ratio is needed.
Third, the judge still mis-scores some short answers, so human calibration is
the most important validation step.

Concretely, we plan to: complete LLM-as-Judge calibration on 20--50 examples and
report Cohen's $\kappa$; tune Provence thresholds and re-run S2/S4; add a
compression-aware coverage metric; scale Notebook~5 from 50 to 200 samples;
prompt the generator for a concise final answer and report EM/F1 before and
after extraction; replace raw surprisal with stronger importance proxies
(question-conditioned similarity, LOO~$\Delta$NLL, learned distillation) and a
ranking/top-$k$ pruning rule before re-running the CfC against full-context,
random, and cosine-history baselines; and migrate the final evaluation from the
HotpotQA pilot to real QReCC dialogue histories.

\begin{credits}
\subsubsection{\discintname}
The authors have no competing interests to declare that are relevant to the
content of this article.
\end{credits}
%
\begin{thebibliography}{99}

\bibitem{lewis2020rag}
Lewis, P., Perez, E., Piktus, A., Petroni, F., Karpukhin, V., Goyal, N.,
K\"{u}ttler, H., Lewis, M., Yih, W., Rockt\"{a}schel, T., Riedel, S., Kiela, D.:
Retrieval-Augmented Generation for Knowledge-Intensive NLP Tasks.
In: Advances in Neural Information Processing Systems (NeurIPS) (2020)

\bibitem{shuster2021retrieval}
Shuster, K., Poff, S., Chen, M., Kiela, D., Weston, J.:
Retrieval Augmentation Reduces Hallucination in Conversation.
In: Findings of EMNLP (2021)

\bibitem{xu2023recomp}
Xu, F., Shi, W., Choi, E.:
RECOMP: Improving Retrieval-Augmented LMs with Compression and Selective Augmentation.
In: International Conference on Learning Representations (ICLR) (2024)

\bibitem{jiang2023llmlingua}
Jiang, H., Wu, Q., Lin, C.-Y., Yang, Y., Qiu, L.:
LLMLingua: Compressing Prompts for Accelerated Inference of Large Language Models.
In: Proceedings of EMNLP (2023)

\bibitem{pan2024llmlingua2}
Pan, Z., Wu, Q., Jiang, H., Garg, M., Liden, N., Ronen, Y., Qiu, L.:
LLMLingua-2: Data Distillation for Efficient and Faithful Task-Agnostic Prompt Compression.
In: Findings of ACL (2024)

\bibitem{yoon2024provence}
Chirkova, N., Yoon, S., Xu, W., Xu, F., Ji, H.:
Provence: Efficient and Robust Context Pruning for Retrieval-Augmented Generation.
In: Proceedings of NAACL (2025)

\bibitem{zheng2023mtbench}
Zheng, L., Chiang, W.-L., Sheng, Y., Zhuang, S., Wu, Z., Zhuang, Y., Lin, Z.,
Li, Z., Li, D., Xing, E.P., Zhang, H., Gonzalez, J.E., Stoica, I.:
Judging LLM-as-a-Judge with MT-Bench and Chatbot Arena.
In: Advances in Neural Information Processing Systems (NeurIPS) (2023)

\bibitem{es2023ragas}
Es, S., James, J., Espinosa-Anke, L., Schockaert, S.:
RAGAS: Automated Evaluation of Retrieval Augmented Generation.
In: Proceedings of EACL (2024)

\bibitem{min2023factscore}
Min, S., Krishna, K., Lyu, X., Lewis, M., Yih, W., Koh, P.W., Iyyer, M.,
Zettlemoyer, L., Hajishirzi, H.:
FActScore: Fine-grained Atomic Evaluation of Factual Precision in Long Form Text Generation.
In: Proceedings of EMNLP (2023)

\bibitem{yang2018hotpotqa}
Yang, Z., Qi, P., Zhang, S., Bengio, Y., Cohen, W.W., Salakhutdinov, R.,
Manning, C.D.:
HotpotQA: A Dataset for Diverse, Explainable Multi-hop Question Answering.
In: Proceedings of EMNLP (2018)

\bibitem{anantha2021qrecc}
Anantha, R., Vakulenko, S., Tu, Z., Longpre, S., Pulman, S., Chappidi, S.:
Open-Domain Question Answering Goes Conversational via Question Rewriting.
In: Proceedings of NAACL (2021)

\bibitem{liu2024lost}
Liu, N.F., Lin, K., Hewitt, J., Paranjape, A., Bevilacqua, M., Petroni, F.,
Liang, P.:
Lost in the Middle: How Language Models Use Long Contexts.
Transactions of the Association for Computational Linguistics (TACL) (2024)

\bibitem{hasani2022cfc}
Hasani, R., Lechner, M., Amini, A., Liebenwein, L., Ray, A., Tschaikowski, M.,
Teschl, G., Rus, D.:
Closed-form Continuous-time Neural Networks.
Nature Machine Intelligence \textbf{4}, 992--1003 (2022)

\bibitem{levy2008}
Levy, R.:
Expectation-based Syntactic Comprehension.
Cognition \textbf{106}(3), 1126--1177 (2008)

\bibitem{smith2013}
Smith, N.J., Levy, R.:
The Effect of Word Predictability on Reading Time is Logarithmic.
Cognition \textbf{128}(3), 302--319 (2013)

\end{thebibliography}
\end{document}
